%#!pdfpLaTeX
%
% 北村研究室用卒論予稿のTeXテンプレートファイル
% 本ファイルは非公式であり，下記で公開されているワードのテンプレートが公式である．
% https://www.kagawa-nct.ac.jp/EE/local/index.html （学内限定アクセス）
%
% 2020年1月18日 北村大地作成
%

\documentclass[a4j]{jsarticle}
\usepackage{nitkagawaproc5ec} % 予稿テンプレートクラスファイル
\usepackage{amsmath,amssymb} % 数式環境
\usepackage{bm} % 数式の太字斜体
\usepackage[dvipdfmx]{color}
\usepackage[dvipdfmx]{graphicx}
\usepackage{flushend} % 2段組みの最終ページの高さを揃える

\renewcommand{\baselinestretch}{0.8} % 行間設定（標準は0.8）

%%%%%%%%%%%%%%%%%%%%%%%%%%% 論文情報 %%%%%%%%%%%%%%%%%%%%%%%%%%%

%%%%% タイトル %%%%%
\title{日本語の研究タイトルをここに記述\\タイトルが2行になる場合はこの行にも記述可能}
\etitle{English title of research theme goes here}

%%%%% 著者 %%%%%
\author{高松 太郎}
\eauthor{Tarou Takamatsu}


%----- 図表題が英語の場合は次の2行を有効化 -----%
%\renewcommand{\figurename}{Fig.~} 
%\renewcommand{\tablename}{Table~}
%--------------------------------------------%

\begin{document}
\maketitle% タイトル生成

%----- ハイフン付きページ番号を表示する場合は次3行を有効化 -----%
%\setcounter{page}{1} % 開始ページ番号（3にすると3ページと4ページの2枚になる）
%\pagestyle{hyphenpage}
%\thispagestyle{hyphenpage}
%---------------------------------------------------------%

%----- ハイフン無しページ番号を表示する場合は次3行を有効化 -----%
%\setcounter{page}{1} % 開始ページ番号（3にすると3ページと4ページの2枚になる）
%\thispagestyle{plain}
%\pagestyle{plain}
%---------------------------------------------------------%

%----- ページ番号を削除する場合は次3行を有効化 -----%
\thispagestyle{empty}
\pagestyle{empty}
%-----------------------------------------------%

%%%%%%%%%%%%%%%%%%%%%%%%%%% 本文 %%%%%%%%%%%%%%%%%%%%%%%%%%%
%----------------------------------------------
\section{はじめに}
%----------------------------------------------

このファイルは卒業論文の予稿のワードテンプレートである．
このテンプレートに記載されている注意事項を熟読したうえ，
予稿の作成に取り掛かること．

%----------------------------------------------
\section{本文の記述方法}
%----------------------------------------------

%----------------------------------------------
\subsection{フォーマット}
%----------------------------------------------

フォーマットはA4縦用紙に2段組み，2ページで作成する．
余白は上下が20~mm，左右が25~mmである．

%----------------------------------------------
\subsection{フォントと各種記号}
%----------------------------------------------

「1. はじめに」等の章タイトルは11~ptのゴシック系フォント
（あるいはそれに準ずるフォント）で記述し，「2.1 フォーマットと
フォント」等の節タイトルは10 ptのゴシック系フォントとする．
本文は全て10~ptの明朝系フォントとし，本文中の半角英数
文字は明朝フォントまたはTimes系フォントとする（本テンプレート
ではTimes系に設定している）．なお，全角の英数文字を
使うことが無いように注意すること．

段落の最初は必ず一字下げを挿入する．句読点は「、」と「。」
の組み合わせか「，」と「．」の組み合わせに統一する．
括弧の記号については，（このように）全角の括弧を用いると
適切なスペースが付与されるため，(このような)半角の括弧よりも
見栄えが良くなる．但し，英文中の括弧では必ず半角括弧と
スペースを組み合わせて記述すること．物理単位については，
4.2~mや5.08~kgのように，半角スペースを空けたうえで付与
することが望ましい．

%----------------------------------------------
\section{図表と数式の記述方法}
%----------------------------------------------

%----------------------------------------------
\subsection{図表の挿入方法}
%----------------------------------------------

図表は必ず本文上部または下部に挿入する．本文の途中に
挿入することの内容に注意せよ．図の挿入例を図\ref{fig:sample}に
示す．このとき，図番号と図タイトルは図の中央下部に配置する．
また，印刷時の品質を向上させるために，図は可能な限り
ラスタ形式ではなくベクタ形式で作成し挿入することを推奨する．
さらに，表の挿入例を表\ref{table:sample}に示す．表の場合
は表番号と表タイトルを表の中央上部に配置する．図番号・
表番号及び図タイトル・表タイトルは本文と同じく10~ptの
明朝系フォントとする．なお，読みやすさの向上ために，図表の
上下には1行分の余白を設けることが望ましい
（図\ref{fig:sample}と表\ref{table:sample}の間や上下を
参照のこと）．

%-%-%-%-%-%-%-%-%
\begin{figure}[t]
  \centering
  \vspace{0pt} % 図上部の余白調整
  \includegraphics[width=0.8\columnwidth]{image.pdf}
  \vspace{0pt} % 図とキャプション間の余白調整
  \caption{図の挿入例}
  \vspace{0pt} % キャプション下部の余白調整
  \label{fig:sample}
\end{figure}
%-%-%-%-%-%-%-%-%

%-%-%-%-%-%-%-%-%
\begin{table}[t]
  \centering
  \caption{表の挿入例}
  \label{table:sample}
  \vspace{0pt} % キャプション下部の余白調整
  \begin{tabular}{|c|c|c|c|} \hline
    項目 & 数値1 & 数値2 & 数値3 \\ \hline
    A & 1 & 5 &  9 \\ \hline
    B & 2 & 6 & 10 \\ \hline
    C & 3 & 7 & 11 \\ \hline
    D & 4 & 8 & 12 \\ \hline
  \end{tabular}
  \vspace{10pt} % 表下部の余白調整
\end{table}
%-%-%-%-%-%-%-%-%

%----------------------------------------------
\subsection{数式の挿入方法}
%----------------------------------------------

本文中に数式を挿入する場合は，次に示すように中央寄せ
すると良い．
\begin{align}
\nonumber  a^2 + b^2 = c^2
\end{align}
あるいは，次式のように左寄せに統一しても良い．
\begin{align}
\nonumber a^2 + b^2 = c^2~~~~~~~~~~~~~~~~~~~~~~~~~~~~~~~~~~~~~~~~~~~~~~~~~
\end{align}
縦に幅を取る数式を挿入すると次のようになる．
\begin{align}
\nonumber  f(x) = a_0 + \sum_{n=1}^{\infty} \left( a_n \cos \frac{ n\pi x }{ L } + b_n\sin \frac{ n\pi x }{ L } \right)
\end{align}
なお，数式中に登場した変数等は全て登場直後に説明されなければならない．その一例を次式に示す．
\begin{align}
\nonumber  x = \frac{ -b \pm \sqrt{ b^2 -4ac } }{ 2a }
\end{align}
ここで，$a$は2乗項の係数，$b$は1乗項の係数，$c$は定数項
をそれぞれ示す．本文全体で未定義の変数等が残っていないか
十分注意すること．

一度登場した数式を後の本文で引用する場合は，次式のように
数式の右側に式番号を付与する．
\begin{align}
  A = \pi r^2 \label{eq:radius}
\end{align}
式番号は括弧書きの数字とし，登場する順番に番号を
付与する．式番号を用いて数式を引用する場合は
「式(\ref{eq:radius})」又は「(\ref{eq:radius})式」と記述する．
どちらを用いても良いが，表記は本文全体で統一されている
必要がある．

文中（インライン）に数式を挿入することもできる．例えば，
$f(x) = 1/2 \sin(\alpha + \beta)$のように記述できる．
ただし，縦に幅を取る数式はインラインで記述することは
避けるべきである．また，分数もスラッシュ記号を用いて
縦の幅を節約する等の工夫をして見やすくする．

%----------------------------------------------
\subsection{参考文献の引用の挿入方法}
%----------------------------------------------

本文の適切な箇所で参考文献を引用する場合は，このように
記述する\cite{sample1}．この引用番号を表す記号は，予稿
の末尾に記載されている参考文献リストと対応している\cite{sample2}．
但し，登場する順に番号を付与する．なお，参考文献として
不適切な文書を引用してはいけない．例えば，信憑性が疑わ
れるインターネット上の記事（ブログ，ウィキペディア，ニュース等）
や学術的でない雑誌は引用を避けるべきである．査読を経て
掲載された原著論文及び国際会議論文が最も適切であるが，
確立された技術を網羅した教科書等を引用する場合もある．
その他，他者の学位論文（修士論文，博士論文等）を引用
することもあるが，信憑性が疑われるものもあるので十分注意
する必要がある．

%----------------------------------------------
\section{まとめ}
%----------------------------------------------

以上の方法で作成すると，正しい体裁の技術文書を作成する
ことができる．より読みやすく正しい技術文書を作成するためには，
詳細を解説した教科書等を参考にすること．また，予稿完成後
も複数人で念入りにチェックを繰り返し，誤字や脱字，不適切
な記述，数式の誤り等が無いように十分注意すること．

%----------------------------------------------
\section{注意事項}
%----------------------------------------------

本テンプレートは元々Wordでのみ配布されている予稿の
テンプレートを北村大地がLaTeXによる組版で可能な限り忠実に
再現したファイルになります．マージンや余白などには細心の
注意を払いましたが，ボールドフォントの違い等の細かい点は
組版ソフトの違いが存在する以上生じてしまいます．
あくまでも非公式のテンプレートとして自己責任で使用してください．

TeXのコンパイラは恐らく最もスタンダードであるpdfpLaTeXを使用
することを想定しています．他のコンパイラでの動作確認はして
おりません．また，今後の対応の予定もありません．もう北村は
疲れました．


%%%%% 参考文献 %%%%%
\begin{thebibliography}{9}% 10以上の文献数であれば99とする

\bibitem{sample1}
  高松太郎，詫間花子，香川一郎,
  ``1つめの参考文献のタイトル,''
  {\it 2020年○○学会講演論文集}, pp.~278--280, 2020.

\bibitem{sample2}
  T.~Takamatsu, H. Takuma, and I. Kagawa, 
  ``Title of reference article,''
  {\it Proc. International Conference on XXX}, pp.~278--280, 2020.


\end{thebibliography}

\end{document}